\subsection{Exterior hilbertian powers}

Let E be a hilbertian space and n a positive integer. For every permutation \(\sigma\in\mathfrak{S}_{n}\) , let \(\varepsilon_{\sigma}\) denote its signature; put \(a_{n}=\frac{1}{n!}\sum_{\sigma\in\mathfrak{S}_{n}}\varepsilon_{\sigma}p_{\sigma}\) in \(\mathcal{L}(\hat{\mathbf{T}}^{n}(\mathrm{E}))\) (V, p.~29). It is immediate that \(a_{n}\) is an orthoprojector, whose image \(\overline{\mathbf{A}_{n}^{\prime}(\mathrm{E})}\) is the closure in \(\hat{\mathbf{T}}^{n}(\mathrm{E})\) of the space \(\mathbf{A}_{n}^{\prime}(\mathrm{E})\) of all antisymmetric tensors of order \(n(\mathrm{A},\mathrm{III},\S7,\mathrm{No}.4)\) . There exists an isomorphism \(\pi_{n}\) from \(\Lambda^{n}(\mathrm{E})\) onto \(\mathbf{A}_{n}^{\prime}(\mathrm{E})\) which is characterized by

\[
\pi_ {n} (x _ {1} \wedge \ldots \wedge x _ {n}) = \mathbf {a} _ {n} (x _ {1} \otimes \ldots \otimes x _ {n}) = \frac {1}{n !} \sum_ {\sigma \in \mathfrak {S} _ {n}} \varepsilon_ {\sigma} x _ {\sigma (1)} \otimes \ldots \otimes x _ {\sigma (n)}\tag{24}
\]

for \(x_{1}, \ldots, x_{n}\) in E. We can now define a Hausdorff prehilbertian space structure on \(\Lambda^{n}(E)\) by putting

\[
\langle u | v \rangle = n! \left\langle \pi_ {n} (u) \mid \pi_ {n} (v) \right\rangle .\tag{25}
\]

More explicitly, we have (compare with formula (30) of A, III, § 11, No.~5).

\[
\langle x _ {1} \wedge \dots \wedge x _ {n} | y _ {1} \wedge \dots \wedge y _ {n} \rangle = \det (\langle x _ {i} | y _ {j} ^ {\prime} \rangle)\tag{26}
\]

for all \(x_{1},\ldots ,x_{n}\) and \(y_{1},\ldots ,y_{n}\) in E.

Let \(\hat{\mathbf{A}}^n (\mathrm{E})\) denote the completion of the prehilbertian space \(\mathbf{A}^n (\mathrm{E})\) , and \(\hat{\mathbf{A}} (\mathrm{E})\) the external hilbertian sum of the hilbertian spaces \(\hat{\mathbf{A}}^n (\mathrm{E})\) .

Example. --- * With the notations of example 1 of V, p.~32, we can identify the hilbertian space \(\symbf{\Lambda}(\mathrm{E})\) with the set of all sequences \((f_n)_{n\geqslant 0}\) of measurable functions for which the number \(\| \mathbf{f}\|\) defined in (22) is finite, and where each function \(f_{n}\) is antisymmetric, that is, satisfies the relation

\[
f _ {n} (\mathbf {x} _ {\sigma (1)}, \dots , \mathbf {x} _ {\sigma (n)}) = \varepsilon_ {\sigma} f _ {n} (\mathbf {x} _ {1}, \dots , \mathbf {x} _ {n})
\]

for every permutation \(\sigma \in \mathfrak{S}_n\) . The hilbertian space \(\hat{\Lambda}(\mathrm{E})\) is called the antisymmetric Fock space corresponding to the weight \(\omega\) . *

PROPOSITION 5. --- Let \((e_i)_{i \in I}\) be an orthonormal basis of the hilbertian space E. Endow I with a totally ordered structure. Then the set of all elements \(e_{i_1} \wedge \ldots \wedge e_{i_n}\) for \(i_1 < \cdots < i_n\) is an orthonormal basis of \(\hat{\Lambda}^n(E)\) .

We know (A, III, § 7, No.~8) that the elements in question form a basis of the vector space \(\Lambda^n (\mathrm{E}_0)\) where \(\mathrm{E}_0\) is the vector subspace of E generated by the vectors \(e_i\) . Further, for \(i_1 < \dots < i_n\) , the matrix of scalar products \(\langle e_{i_k}|e_{i_\ell}\rangle\) is the unit matrix of order \(n\) ; by (26), we have \(\| e_{i_1}\wedge \ldots \wedge e_{i_n}\| = 1\) . Finally, if \((i_1,\dots,i_n)\) and \((j_1,\dots,j_n)\) are two distinct, strictly increasing sequences of elements of I, then there exists an element \(j_\ell\) distinct from \(i_1,\ldots ,i_n\) and so we have \(\langle e_{i_k}|e_{j_\ell}\rangle = 0\) for \(1\leqslant k\leqslant n\) , and by (26), \(\langle e_{i_1}\wedge \ldots \wedge e_{i_n}|e_{j_1}\wedge \ldots \wedge e_{j_n}\rangle = 0\) . In other words, the family of elements \(e_{i_1}\wedge \ldots \wedge e_{i_n}\) , for \(i_1 < \dots < i_n\) , is orthonormal.

But \(E_{0}\) is dense in E, and the mapping \((x_{1},...,x_{n})\mapsto x_{1}\wedge\ldots\wedge x_{n}\) from \(E\times\cdots\times E\) into \(\Lambda^{n}(E)\) is continuous. Consequently, \(\Lambda^{n}(E_{0})\) is dense in \(\Lambda^{n}(E)\) , and proposition 5 follows.

COROLLARY. --- Suppose that the hilbertian space E is the direct sum of two orthogonal subspaces M and N. The canonical isomorphism g from \(\symbf{\Lambda}(\mathbf{M})\otimes\symbf{\Lambda}(\mathbf{N})\) onto \(\symbf{\Lambda}(\mathbf{E})\) (A, III, § 7, No.~7) extends in a unique way to a hilbertian space isomorphism from \(\hat{\symbf{\Lambda}}(\mathbf{M})\hat{\otimes}_{2}\hat{\symbf{\Lambda}}(\mathbf{N})\) onto \(\hat{\symbf{\Lambda}}(\mathbf{E})\) .

The proof is analogous to that of the corollary of prop. 4 (V, p.~31).

Let \(\mathbf{E}\) and \(\mathbf{F}\) be two hilbertian spaces and \(u \in \mathcal{L}(\mathbf{E};\mathbf{F})\) . We shall show, as in the case of symmetric powers \(\hat{\mathbf{S}}^n (\mathbf{E})\) (V, p.~32) that the linear mapping \(\symbf{\Lambda}^n (u)\) from \(\symbf{\Lambda}^n (\mathbf{E})\) into \(\symbf{\Lambda}^n (\mathbf{F})(\mathbf{A},\mathbf{III},\S 7,\mathbf{No}.4)\) extends to a continuous linear mapping \(\hat{\symbf{\Lambda}}^n (u)\) from \(\hat{\symbf{\Lambda}}^n (\mathbf{E})\) into \(\hat{\symbf{\Lambda}}^n (\mathbf{F})\) . We have the relations

\[
\hat {\symbf {\Lambda}} ^ {n} (1 _ {\mathrm{E}}) = 1 _ {\hat {\symbf {\Lambda}} ^ {n} (\mathrm{E})},\tag{27}
\]

\[
\hat {\symbf {\Lambda}} ^ {n} (v \circ u) = \hat {\symbf {\Lambda}} ^ {n} (v) \circ \hat {\symbf {\Lambda}} ^ {n} (u) \quad \text { if } v \text { belongs   to } \mathscr {L} (\mathrm{F}; \mathrm{G}),\tag{28}
\]

\[
\left\| \hat {\symbf {\Lambda}} ^ {n} (u) \right\| \leqslant \| u \| ^ {n}.\tag{29}
\]

In general we do not have equality in formula (29) (TS, IV, § 6). Finally, we have an isomorphism \(\psi_{n} = \psi_{n,\mathrm{E}}\) from \(\hat{\mathbf{A}}^{n}(\mathrm{E})\) onto the subspace \(\overline{\mathbf{A}_{n}^{\prime}(\mathrm{E})}\) of \(\hat{\mathbf{T}}^{n}(\mathrm{E})\) defined by

\[
\psi_ {n} (x _ {1} \wedge \dots \wedge x _ {n}) = \frac {1}{(n !) ^ {1 / 2}} \sum_ {\sigma \in \mathfrak {S} _ {n}} \varepsilon_ {\sigma} x _ {\sigma (1)} \otimes \dots \otimes x _ {\sigma (n)}.\tag{30}
\]

